\centerline{\bf \Huge Круги Эйлера I}


\begin{flushright}
        \scriptsize{20:31 прибыл Годжо Сатору... и решил все задачи}
\end{flushright}

\problem{\textbf{1}} В классе все увлекаются математикой или биологией. Сколько
человек в классе, если математикой занимаются 15 человек, биологией — 20, а математикой
и биологией — 10?


\problem{\textbf{2}} На каникулах в кино пришло 200 ребят. На приключенческий фильм было продано 152 билета, а на комедию — 97.
Сколько ребят посмотрели и тот фильм, и другой? (Каждый посмотрел по меньшей мере один из фильмов.)


\problem{\textbf{3}} В кондитерском отделе супермаркета посетители обычно покупают либо один торт, либо одну коробку конфет, либо один торт и одну коробку конфет. В один из дней было продано 57 тортов и 36 коробок конфет. Сколько было покупателей, если 12 человек купили и торт, и коробку конфет?

\begin{flushright} \textbf{2 балла}\\ 
\end{flushright}

\problem{\textbf{4}} Во дворе стоят машины. Некоторые из них — москвичи, а остальные — жигули. Некоторые из машин красные, а остальные белые. Некоторые из машин новые, а остальные — старые. Известно, что красных москвичей — 3, новых москвичей — 4, а новых красных машин — 5. При этом старых белых москвичей — 2, новых белых жигулей — 1, а старых красных москвичей вообще ни одного. Сколько во дворе новых красных москвичей, если всего машин 21, а старых белых жигулей — 6?

\problem{\textbf{5}} Из 100 ребят, отправляющихся в детский оздоровительный лагерь, кататься на сноуборде умеют 30 ребят, на скейтборде — 28, на роликах — 42. На скейтборде и на сноуборде умеют кататься 8 ребят, на скейтборде и на роликах — 10, на сноуборде и на роликах — 5, а на всех трех — 3. Сколько ребят не умеют кататься ни на сноуборде, ни на скейтборде, ни на роликах? (В число умеющих кататься на сноуборде включены те, кто умеет кататься ещё на чём-либо, и так далее).

\begin{flushright} \textbf{4 балла}\\ 
\end{flushright}

